%! Representing a Categorical Variable — Unit 1, Lesson 1.2 — Experience & Formalize
%! (Activity · QuickNotes · Application)
% EFFL component, Math Medic sizing: 12pt body, OPEN white space after every question.
% \answerspace{H}{} reserves exactly H in the blank; the key passes the answer as the second
% argument, so the two files paginate identically by construction.
% Page budget: Activity <=2pp · QuickNotes 1/2pp · Application 1/2-1pp.
% Check Your Understanding is NOT in this component: those items were moved to the end of
% the packet and are now the lesson's SCORED homework (see ../homework/).
%
% SPOILER RULE: the Activity never says frequency table, relative frequency, proportion, or
% bar graph. Column heads say "Percent"; the pre-drawn displays are called "the chart the
% principal posted". `Categorical'/`quantitative' WERE formalized in Lesson 1.1, so they are
% prior knowledge here, not spoilers.
%
% NO SKETCH QUESTIONS (hard constraint): every display is pre-drawn in TikZ and read or
% critiqued. Students "construct" only by filling in a table, never by drawing an axis.
%
% The \begin{work} block in scenario 2b is authored BYTE-IDENTICALLY here and in
% experience_key/ — that block plus \answerspace is what keeps the two page-for-page.
\documentclass[12pt]{article}
\usepackage{apstats-article}
\usepackage{apstats-boxes}
\newcommand{\answerspace}[2]{\par\nopagebreak\noindent\begin{minipage}[t][#1][t]{\linewidth}%
  \color{keyred}\bfseries #2\end{minipage}\par}

\begin{document}

\pageheader{Unit 1, Lesson 1.2}{Experience \& Formalize: Getting to School}
% NAMESTRIP: no \namedateperiod here — the name row lives on the cover only.

\vspace{0.02in}

% ── Part 1: the Activity ─────────────────────────────────────────────────────
% Scenario 1 builds the table toolkit on ONE school and plants the plurality-vs-majority
% trap in 1d. Scenario 2 carries the CRUX: South has more walkers AND a lower walking rate,
% so "which school walks more" has two defensible answers depending on the number you use.
\begin{headlinebox}{sky}
\textbf{Getting to School.} A district asked every student at two of its high schools one
question: \emph{how do you usually get to school?} The results are below. Work through the two
problems with your group using only what you already know, and be ready to say \emph{which
number} you used for each answer.
\end{headlinebox}

\vspace{0.04in}

\begin{scenariobox}[1. One School]{navy}
North High surveyed all \textbf{400} of its students.

\vspace{5pt}
\normalsize
\begin{enumerate}[label=\textbf{\alph*.}, leftmargin=1.8em, itemsep=6pt]
  \item Complete the last column.

        \vspace{3pt}
        \renewcommand{\arraystretch}{1.35}
        \begin{tabularx}{\linewidth}{@{} l Y Y @{}}
        \rowcolor{sky}
        \textbf{How they get to school} & \textbf{Number of students} &
        \textbf{Percent of all 400} \\
        \hline
        Bus  & 180 & \blank{2.4cm} \\
        Car  & 120 & \blank{2.4cm} \\
        Walk &  80 & \blank{2.4cm} \\
        Bike &  20 & \blank{2.4cm} \\
        \hline
        \textbf{Total} & \blank{2.0cm} & \blank{2.4cm} \\
        \end{tabularx}
  \item The two totals at the bottom are not the same kind of thing. Say what each one has to
        come out to, and why it has to.
        \answerspace{1.6cm}{}
  \item A student council flyer says \textbf{``Most North students take the bus.''} Use your
        table to decide whether that is a fair thing to print, and say what ``most'' would
        require.
        \answerspace{1.8cm}{}
\end{enumerate}
\end{scenariobox}

\boxguard[16]
\begin{scenariobox}[2. Two Schools]{navy}
South High is bigger: it surveyed all \textbf{1{,}000} of its students.

\vspace{5pt}
\small
\renewcommand{\arraystretch}{1.2}
\begin{tabularx}{\linewidth}{@{} l Y Y Y Y Y @{}}
\rowcolor{sky}
\textbf{School} & \textbf{Bus} & \textbf{Car} & \textbf{Walk} & \textbf{Bike} &
\textbf{Total surveyed} \\
\hline
North & 180 & 120 &  80 & 20 &   400 \\
South & 500 & 330 & 150 & 20 & 1{,}000 \\
\hline
\end{tabularx}

\normalsize
\vspace{4pt}
The district posted this chart to compare the two schools.
\vspace{4pt}
\begin{center}
\begin{tikzpicture}[scale=0.85]
  \draw[linegray, very thin] (0,1) -- (6.6,1);
  \draw[linegray, very thin] (0,2) -- (6.6,2);
  \draw[linegray, very thin] (0,3) -- (6.6,3);
  \draw[->, thick] (0,0) -- (0,3.8);
  \draw[->, thick] (0,0) -- (6.8,0);
  \foreach \y/\lab in {1/200, 2/400, 3/600} {
    \draw (-0.08,\y) -- (0.08,\y) node[left, font=\scriptsize, xshift=-3pt] {\lab};}
  \foreach \x/\hn/\hs/\lab in {0.4/0.9/2.5/Bus, 2.0/0.6/1.65/Car, 3.6/0.4/0.75/Walk,
                               5.2/0.1/0.1/Bike} {
    \fill[navy]   (\x,0)     rectangle ++(0.55,\hn);
    \fill[goldacc](\x+0.6,0) rectangle ++(0.55,\hs);
    \node[below, font=\footnotesize] at (\x+0.57,-0.05) {\lab};}
  \node[rotate=90, font=\scriptsize] at (-0.85,1.9) {Students};
  \fill[navy] (5.9,3.3) rectangle ++(0.28,0.28);
  \node[right, font=\scriptsize] at (6.2,3.44) {North};
  \fill[goldacc] (5.9,2.9) rectangle ++(0.28,0.28);
  \node[right, font=\scriptsize] at (6.2,3.04) {South};
\end{tikzpicture}
\end{center}
\vspace{4pt}
\begin{enumerate}[label=\textbf{\alph*.}, leftmargin=1.8em, itemsep=6pt]
  \item Which school has \textbf{more walkers}? \blank{3.0cm} \quad Now work out what
        \emph{percent} of each school walks.
        \begin{work}
        \text{North} &= \frac{80}{400} = 0.20 = 20\% \\[4pt]
        \text{South} &= \frac{150}{1000} = 0.15 = 15\%
        \end{work}
        Which school has the \textbf{higher walking rate}? \blank{3.2cm}
  \item \textbf{Those are two different answers to ``which school walks more.''} Explain how
        both can be true at once. Then look again at the chart: South's bar is taller in
        \emph{every} category, which was guaranteed before anyone counted. Name the one change
        that would make that chart a fair comparison.
        \answerspace{2.8cm}{}
  \item Complete South's percentages, then find every category where the chart's ranking
        \textbf{reverses} once you use percents instead of counts.

        \vspace{2pt}
        \textbf{Bus} \blank{1.7cm} \quad \textbf{Car} \blank{1.7cm} \quad
        \textbf{Walk} \blank{1.7cm} \quad \textbf{Bike} \blank{1.7cm}
        \answerspace{1.8cm}{}
\end{enumerate}
\end{scenariobox}

% ── Part 2: QuickNotes (the debrief fills this) — target: half a page ────────
\boxguard[14]
\begin{tcolorbox}[enhanced, colback=sky, colframe=navy, boxrule=0.8pt, arc=2mm,
    left=3mm, right=3mm, top=1.5mm, bottom=1.5mm, breakable,
    fonttitle=\bfseries\small\color{white}, title={QuickNotes: Tables and Graphs for a Categorical Variable},
    attach boxed title to top left={yshift=-2mm, xshift=4mm},
    boxed title style={colback=navy, colframe=navy, arc=1.5mm}]
\small
\begin{itemize}[leftmargin=1.1em, itemsep=4pt]
  \item A \textbf{frequency table} gives the \blank{2.0cm} of cases in each category. A
        \textbf{relative frequency table} gives the \blank{2.4cm} in each category.
  \item Frequencies add to the \blank{2.6cm}; relative frequencies always add to
        \blank{1.4cm} (or 100\%).
  \item Percents, proportions, and rates all carry the \blank{2.0cm} information --- pick
        whichever is easiest to read.
  \item In a \textbf{bar graph}, the height of each bar is the \blank{1.8cm} or the
        \blank{2.4cm} for that category. Bars are drawn with \blank{2.0cm} between them,
        because the categories are separate groups and not a number line.
  \item \textbf{Comparing two groups of different sizes: always use \blank{3.4cm}, never raw
        counts.} The bigger group wins on counts before anyone counts anything.
  \item ``Most'' means more than \blank{1.6cm} of the total. The biggest category is only the
        \emph{largest share} --- it need not be a majority.
\end{itemize}
\end{tcolorbox}

% ── Part 3: the Application (worked together, ~6 min) ────────────────────────
\boxguard[14]
\begin{notesbox}{Application: Two Shifts}
We will do this one together. A coffee shop records the size of every drink it sells. The
morning shift filled \textbf{240} orders; the evening shift filled \textbf{80}.

\vspace{4pt}
\small
\renewcommand{\arraystretch}{1.2}
\begin{tabularx}{\linewidth}{@{} l Y Y Y Y @{}}
\rowcolor{sky}
\textbf{Shift} & \textbf{Small} & \textbf{Medium} & \textbf{Large} & \textbf{Total} \\
\hline
Morning &  48 & 108 & 84 & 240 \\
Evening &  18 &  30 & 32 &  80 \\
\hline
\end{tabularx}

\normalsize
\begin{enumerate}[label=\textbf{\alph*.}, leftmargin=1.8em, itemsep=6pt]
  \item Which shift sold more large drinks? \blank{2.6cm} \quad
        Which shift sells large drinks at the higher rate? \blank{2.6cm}
  \item The owner wants one sentence for the staff newsletter comparing what the two shifts
        sell. Write it, and say which kind of number you used and why.
        \answerspace{1.6cm}{}
  \item \textbf{What if we change the question?} The owner now asks, ``Which shift is busier?''
        Which column of this table answers that --- and is it a rate or a count?
        \answerspace{1.4cm}{}
\end{enumerate}
\end{notesbox}

\end{document}
